\documentclass[11pt,a4paper]{article}
\usepackage[T1]{fontenc}
\usepackage[utf8]{inputenc}
\usepackage[margin=26mm,headheight=14pt]{geometry}
\usepackage{amsmath,amsthm,mathtools}
\usepackage{newtxtext,newtxmath}
\usepackage{microtype,booktabs,array,longtable}
\usepackage[dvipsnames]{xcolor}
\usepackage{enumitem,fancyhdr,xurl}
\usepackage[colorlinks=true,linkcolor=MidnightBlue,citecolor=MidnightBlue,
 urlcolor=MidnightBlue,pdftitle={Sharp Fixed-Phase Order of Rvachev Quadrature},
 pdfauthor={Research manuscript prepared with ChatGPT},
 pdfsubject={Dyadic maximality, effective refinement thresholds, and defect geometry}]{hyperref}
\usepackage[capitalise,noabbrev]{cleveref}
\numberwithin{equation}{section}
\setlength{\parskip}{3pt}
\setlength{\parindent}{1.2em}
\setlist{itemsep=3pt,topsep=5pt}
\pagestyle{fancy}
\fancyhf{}
\fancyhead[L]{\small Sharp fixed-phase order of Rvachev quadrature}
\fancyhead[R]{\small Research manuscript}
\fancyfoot[C]{\thepage}
\renewcommand{\headrulewidth}{0.3pt}
\theoremstyle{plain}
\newtheorem{theorem}{Theorem}[section]
\newtheorem{lemma}[theorem]{Lemma}
\newtheorem{proposition}[theorem]{Proposition}
\newtheorem{corollary}[theorem]{Corollary}
\theoremstyle{definition}
\newtheorem{definition}[theorem]{Definition}
\newtheorem{example}[theorem]{Example}
\theoremstyle{remark}
\newtheorem{remark}[theorem]{Remark}
% ed. (2026-09-29): unnumbered environment for editorial notes.
\newtheorem*{ednote}{Editorial note (ProveIt, 2026-09-29)}
\newcommand{\R}{\mathbb R}
\newcommand{\C}{\mathbb C}
\newcommand{\N}{\mathbb N}
\newcommand{\Z}{\mathbb Z}
\newcommand{\T}{\mathbb R/\mathbb Z}
\newcommand{\dd}{\,\mathrm d}
\newcommand{\ii}{\mathrm i}
\newcommand{\sinc}{\operatorname{sinc}}
\newcommand{\supp}{\operatorname{supp}}
\newcommand{\dist}{\operatorname{dist}}
\newcommand{\sgn}{\operatorname{sgn}}
\newcommand{\ord}{\operatorname{ord}}
\newcommand{\vTwo}{\nu_2}
\newcommand{\Poly}{\mathcal P}
\newcommand{\norm}[1]{\left\lVert#1\right\rVert}
\newcommand{\abs}[1]{\left|#1\right|}
\newcommand{\code}[1]{\texttt{\small #1}}
\newcommand{\oddSum}{\sum_{\substack{n\ge1\\n\text{ odd}}}}

\begin{document}
% ed. (2026-09-29): no page anchor on the title page (it duplicated page.1).
\hypersetup{pageanchor=false}
\begin{titlepage}
\centering
\vspace*{10mm}
{\small\scshape Analysis / Fabius function / Rvachev up-function\par}
\vspace{14mm}
{\LARGE\bfseries Sharp Fixed-Phase Order\par}
\vspace{3mm}
{\LARGE\bfseries of Rvachev Quadrature\par}
\vspace{9mm}
{\large Dyadic maximality, effective refinement thresholds,\par}
\vspace{2mm}
{\large and quantitative defect geometry\par}
\vspace{12mm}
{Research manuscript prepared with ChatGPT\par}
{29 September 2026\par}
\vspace{12mm}
\begin{minipage}{0.93\textwidth}
\begin{abstract}
A recent ProveIt manuscript classifies the phases at which the Rvachev
lattice rule gains one polynomial degree, but leaves open whether a selected
phase can gain a second degree. We settle this maximality question at every
dyadic mesh. If $M=2^d$, the largest degree attainable at any fixed phase is
exactly $d+1$. The two phases attaining it both fail at degree $d+2$, with an
explicit signed leading term and a relative error smaller than $1/4$.

For $M=2^{r-1}m$, with $m$ odd, we derive an exact next-defect formula and a
computable sufficient condition for the same maximality conclusion. It
holds eventually for every fixed $m$. More explicitly, if
$J=\max\{2,1+\lceil\log_2m\rceil\}$, it holds for every
$r\ge2J(m+1)+4$. Thus any exception at a fixed odd part lies in an explicitly
bounded finite set of refinement levels. A refined expansion separates the
$2^{-r}$ and $3^{-r}$ corrections and explains their parity dependence.

A three-factor Fourier estimate also shows that the first defect has
exactly two nondegenerate stationary points, and gives a quantitative
stability theorem for its zeros under $C^1$ perturbations. The results
transfer to moment-deconvolved polynomial synthesis and to refined phase
averages. The article supplies complete analytic proofs, independent exact
and numerical checks, and a proposed formalization and research program.
\end{abstract}
\end{minipage}
\vfill
\begin{minipage}{0.93\textwidth}
\small
\textbf{Status.} These are conventional mathematical proofs, not new
Lean-verified theorems. The general all-integer-mesh maximality question is
not claimed solved. Novelty is assessed relative to the inspected repository
sources and primary references, not certified against the entire literature.
This is an unrefereed research manuscript.
\end{minipage}
\end{titlepage}
% ed. (2026-09-29): page anchors resume after the title page.
\hypersetup{pageanchor=true}
\begingroup
\small\setlength{\parskip}{0pt}
\tableofcontents
\endgroup
\newpage

\section{The research question and the contribution}
\label{sec:intro}

\subsection{From phase completeness to fixed-phase maximality}

Let $u$ be the normalized Rvachev up-function: a nonnegative even
$C^\infty$ function supported on $[-1,1]$, with integral one. Its Fourier
transform, with the convention fixed below, is
\begin{equation}\label{eq:H}
 H(t)=\widehat u(t)=\prod_{j=0}^{\infty}
 \sinc\!\left(\frac{\pi t}{2^j}\right),
 \qquad \sinc z=\begin{cases}\sin z/z,&z\ne0,\\1,&z=0.\end{cases}
\end{equation}
For $M>0$ and a phase $\theta\in\T$, define
\begin{equation}\label{eq:rule}
 Q_{M,\theta}(P)=\frac1M\sum_{k\in\Z}
 P\!\left(\frac{k+\theta}{M}\right)
 u\!\left(\frac{k+\theta}{M}\right),
 \qquad I(P)=\int_\R P(x)u(x)\dd x.
\end{equation}
The actual spacing is $1/M$. The sum is finite. Write $\Poly_s$ for the
polynomials of degree at most $s$, including zero.

For integer $M$, put $d=\vTwo(M)$, the exponent of two dividing $M$.
The repository's module \code{RvachevSuperconvergentSynthesis.lean}
establishes exactness on $\Poly_d$ at every phase, and sufficient special
phases for exactness on $\Poly_{d+1}$ \cite{repo-lean}. The subsequent
manuscript \emph{Complete Superconvergence Phases and Optimal Phase Filters
for Rvachev Quadrature}, dated 28 September 2026, proves that these are all
such phases \cite{repo-phase}. Specifically, with $r=d+1$,
\begin{equation}\label{eq:Z}
 Z_r=\begin{cases}
 \{0,\tfrac12\}\pmod1,&r\text{ odd},\\[2pt]
 \{\tfrac14,\tfrac34\}\pmod1,&r\text{ even}.
 \end{cases}
\end{equation}
The same manuscript explicitly leaves the next-degree failure problem
unresolved. Its results about optimal \emph{uniformly translated} phase
filters do not answer that different question.

% ed. (2026-09-29): relation to the question answered; uncited repository results.
\begin{ednote}
This article answers Question~1 (``Fixed-phase maximality beyond the first
extra degree'') of \cite{repo-phase} for every dyadic mesh and, from the
explicit level of \cref{thm:elementary} on, for every odd part; the finitely
many lower levels of each non-dyadic mesh remain open. An editorial note at
that question records this. The completeness of the phase set \eqref{eq:Z}
at the first failing level is also Theorem \texttt{thm:phase-zero-set} of
the canonical synthesis
\path{inverse-and-sampling/comb-interpolation/comb_interpolation_synthesis/},
and the phase-uniform sharpness of exactness through degree $\vTwo(M)$ is
machine-checked as
\path{Fabius.exists_shift_tsum_shifted_monomial_ne_integral_nat_real}
(\path{Analysis/FabiusFunction/Lean/FabiusFunction/CompositeMeshSharpness.lean});
neither is cited in this article. The fixed-phase results of this article
have no Lean statement.
\end{ednote}

We ask whether, at a selected phase, the degree-$r+1$ monomial still has a
nonzero error. An affirmative answer at a mesh $M$ completes its fixed-phase
classification: all phases outside $Z_r$ have exact degree $r-1$, and the
two phases in $Z_r$ have exact degree $r$.

\subsection{Main conclusions}

The central result is \cref{thm:dyadic}: for every $d\ge0$, the optimal
fixed-phase degree at $M=2^d$ is exactly $d+1$. Its proof is quantitative,
with a uniform relative margin; it is not an extrapolation from small
meshes. The analytic mechanism is the explicit separation of the next
Fourier defect into odd frequencies and a refined-mesh defect.

For a general odd part $m$, \cref{thm:criterion} gives an explicit sufficient
inequality, \cref{thm:eventual} gives a cutoff in terms of two computable
Fourier constants, and \cref{thm:elementary} gives the fully elementary bound
\begin{equation}\label{eq:elementary-intro}
 r\ge2\max\{2,1+\lceil\log_2m\rceil\}(m+1)+4.
\end{equation}
No unproved nonvanishing assumption occurs in that result. The constants
are conservative, but the conclusion reduces each fixed-odd-part problem
to finitely many levels, with an $O(m\log(m+1))$ explicit cutoff.

Two related conclusions are proved along the way. A strengthened
odd-multiplier estimate gives the complete stationary-point geometry of
the first defect and robust zero localization. The next-defect identity
gives a two-exponential asymptotic expansion, with an entire even-frequency
contribution disappearing at one of the two parities.

\subsection{Scope, provenance, and limits}

The repository snapshot inspected for the relevant source files was
\begin{center}\small
\texttt{23dd71d2d3d5e85d2b97a8cc45f55e735d69a4ae}.
\end{center}
The pre-existing phase-completeness theorem, its two-factor estimate, the
uniform filter classification, and classical dyadic evaluations are credited
rather than presented as new. Primary descriptions of $u$ and its arithmetic
are due to Arias de Reyna \cite{arias1982,arias2017}. This article re-proves
the few foundational facts needed to keep the new argument self-contained.

We have not audited the whole repository or independently certified the
entire literature. In particular, the term ``maximality'' here concerns the
fixed-node, fixed-weight family \eqref{eq:rule}; it is not a statement about
all Gaussian quadrature rules or arbitrary weights. We do not resolve every
mesh $2^{r-1}m$ at the finitely many levels below our sufficient thresholds.
The later research program states this remaining gap explicitly.

\section{Fourier normalization and the first nonzero defect}
\label{sec:fourier}

\subsection{Construction and zero multiplicities}

Take independent variables $V_j$, each uniform on $[-1,1]$, and let
\[
 Y=\sum_{j=1}^{\infty}2^{-j}V_j.
\]
The sum is absolutely convergent and belongs to $[-1,1]$. Its characteristic
function in the convention
\[
 \widehat f(t)=\int_\R f(x)e^{-2\pi\ii tx}\dd x
\]
is \eqref{eq:H}. The product converges locally uniformly on $\C$, since
its tail factors differ from one by $O(4^{-j})$ on compact sets. It is an
entire function.

For every fixed $N$, bounding the first $N$ factors by reciprocal absolute
arguments and the remaining real factors by one gives
$H(t)=O_N(|t|^{-N})$. Fourier inversion therefore gives a smooth density
$u$ of $Y$. It is compactly supported, even, nonnegative, and normalized.
Consequently $H$ and all its real derivatives are Schwartz functions. This
also justifies the classical normalization used in \cite{arias1982}.

\begin{lemma}[Real zeros]\label[lemma]{lem:zeros}
The real zeros of $H$ are exactly $\Z\setminus\{0\}$. Their orders are
\begin{equation}\label{eq:order}
 \ord_NH=\vTwo(|N|)+1\qquad(N\in\Z\setminus\{0\}).
\end{equation}
In particular $H(m/2)\ne0$ for odd $m$.
\end{lemma}
\begin{proof}
At a nonzero integer $N$ of valuation $d$, precisely the factors with
$j=0,\ldots,d$ vanish, each to first order. All the remaining factors are
nonzero. Their tail product is nonzero because the sum of the absolute
deviations from one converges. If $t$ is not a nonzero integer, no factor
vanishes, and the same tail argument excludes an additional product zero.
\end{proof}

\subsection{Poisson summation in physical coordinates}

Set
\begin{equation}\label{eq:defect}
 E_{s,M}(\theta)=Q_{M,\theta}(x^s)-I(x^s).
\end{equation}
The periodization of the Schwartz function $x^su(x)$ gives
\begin{equation}\label{eq:poisson}
 E_{s,M}(\theta)=\left(\frac{\ii}{2\pi}\right)^s
 \sum_{n\in\Z\setminus\{0\}}H^{(s)}(Mn)e^{2\pi\ii n\theta}.
\end{equation}
For completeness, the $n$th Fourier coefficient of
$M^{-1}\sum_k f((k+\theta)/M)$ is $\widehat f(Mn)$, by the substitution
$x=(k+\theta)/M$. Differentiation under the integral gives
$\widehat{x^su}(t)=(\ii/(2\pi))^sH^{(s)}(t)$. Removing the constant
coefficient proves \eqref{eq:poisson}. All phase derivatives of the series
converge absolutely and uniformly.

From now on write
\begin{equation}\label{eq:mesh}
 M=2^{r-1}m,\qquad r\ge1,\qquad m\ge1\text{ odd}.
\end{equation}
By \cref{lem:zeros}, $E_{s,M}\equiv0$ for $s<r$. Thus only the degree-$r$
and degree-$(r+1)$ defects are needed for a sharp classification.

\subsection{Normalized coefficients}

Define the nonzero real number and real coefficients
\begin{equation}\label{eq:ab}
 A_m=H(m/2),\qquad
 a_{m,n}=\frac{H(mn/2)}{A_m},\qquad
 b_{m,n}=\frac{m}{2}\frac{H'(mn/2)}{A_m}
 \quad(n\ge1\text{ odd}),
\end{equation}
so $a_{m,1}=1$. We often omit the subscript $m$ inside proofs. Put
\begin{equation}\label{eq:beta}
 \beta_m=-b_{m,1}=-\frac m2\frac{H'(m/2)}{H(m/2)},
 \qquad \tau_s(\theta)=\cos(2\pi\theta+s\pi/2).
\end{equation}
The letter $\mathcal F$ below denotes a normalized defect, not the Fabius
function:
\begin{equation}\label{eq:F}
 \mathcal F_{s,m}(\theta)=
 \oddSum \frac{a_{m,n}}{n^s}\tau_s(n\theta),\qquad s\ge1.
\end{equation}
Its actual coefficients are rapidly decreasing for fixed $m$, so the
series is smooth.

\begin{proposition}[First defect; established phase framework]
\label[proposition]{prop:first}
For the mesh \eqref{eq:mesh},
\begin{equation}\label{eq:first}
 E_{r,M}(\theta)=
 -\frac{2r!A_m}{(2\pi M)^r}\mathcal F_{r,m}(\theta).
\end{equation}
\end{proposition}
\begin{proof}
At $N=Mn$ with $n$ odd, split the product after its first $r$ factors.
Exactly these factors vanish, and their first-order product is
$-(z-N)^r/N^r$: among their integer arguments $N/2^j$, exactly the last
is odd. Hence
\[
 H^{(r)}(Mn)=-\frac{r!}{(Mn)^r}H(mn/2).
\]
For even $n$, the derivative vanishes. Pair positive and negative odd
frequencies in \eqref{eq:poisson}, using
$H^{(r)}(-t)=(-1)^rH^{(r)}(t)$, to obtain \eqref{eq:first}.
\end{proof}

\section{A strengthened multiplier estimate and defect geometry}
\label{sec:geometry}

Write
\begin{equation}\label{eq:S}
 S_s=\sum_{\substack{n\ge3\\n\text{ odd}}}n^{-s}
 =(1-2^{-s})\zeta(s)-1\qquad(s>1).
\end{equation}
The following elementary bound will be useful:
\begin{equation}\label{eq:Sbound}
 S_s\le3^{-s}+\frac12\int_3^\infty x^{-s}\dd x
 =3^{-s}\left(1+\frac{3}{2(s-1)}\right).
\end{equation}
Indeed, for every odd $n\ge5$, the integral over $[n-2,n]$ is at least
$2n^{-s}$.

\begin{lemma}[Two-factor and three-factor domination]\label[lemma]{lem:multiplier}
For positive odd $m,n$,
\begin{equation}\label{eq:multiplier}
 |a_{m,n}|\le\min\left\{\frac1{n^2},
 \frac{1+\sqrt2}{n^3}\right\}.
\end{equation}
\end{lemma}
\begin{proof}
For integer $n\ge1$ and real $x$, $|\sin(nx)|\le n|\sin x|$.
One proof expands $\sin(nx)/\sin x$ as a sum of $n$ complex numbers
of modulus one when the denominator is nonzero, and then uses continuity.
Thus every individual sinc-factor ratio in $H(mn/2)/H(m/2)$ has
absolute value at most one.

For $j=0$, the arguments are odd multiples of $\pi/2$, so the sine
magnitudes are both one and the ratio is exactly $1/n$. For $j=1$,
the sine magnitudes at odd multiples of $\pi/4$ are both $1/\sqrt2$,
so that ratio is also $1/n$. This proves the first bound, the estimate
used in the earlier phase-completeness proof \cite{repo-phase}.

At $j=2$, the sine magnitudes at odd multiples of $\pi/8$ belong to
$\{\sin(\pi/8),\sin(3\pi/8)\}$. Their ratio is at most
$\sin(3\pi/8)/\sin(\pi/8)=1+\sqrt2$. That factor contributes at
most $(1+\sqrt2)/n$, and all later ratios are at most one. Passing from
finite products to their nonzero limits proves the second bound.
\end{proof}

\begin{theorem}[Two zeros and exactly two stationary points]
\label{thm:shape}
Put $c_r=S_{r+1}$ and $\eta_r=(1+\sqrt2)S_{r+1}$. Then, for every odd
$m$ and $r\ge1$,
\begin{align}
 |\mathcal F_{r,m}(\theta)-\tau_r(\theta)|
 &\le c_r|\tau_r(\theta)|,\label{eq:C0}\\
 |\mathcal F'_{r,m}(\theta)-\tau'_r(\theta)|
 &\le\eta_r|\tau'_r(\theta)|.\label{eq:C1}
\end{align}
Here $c_r\le\pi^2/8-1<1/4$ and $\eta_r<1$. The zeros of
$\mathcal F_{r,m}$ are exactly $Z_r$ and are simple. Its stationary
points are exactly the zeros of $\tau'_r$, and are nondegenerate. It
has the same sign and the same monotonicity intervals as $\tau_r$.
\end{theorem}
\begin{proof}
For odd $n$, the integer multiple inequality, applied to either sine or
cosine according to the parity of $r$, gives
\[
 |\tau_r(n\theta)|\le n|\tau_r(\theta)|.
\]
Combining this with the first bound in \eqref{eq:multiplier} proves
\eqref{eq:C0}. This is the known phase-completeness argument, reproduced
here to establish the exact baseline.

Differentiation contributes a factor $n$, and the same trigonometric
inequality applied to the differentiated harmonic gives
\[
 \abs{\frac{\dd}{\dd\theta}\tau_r(n\theta)}
 \le n^2|\tau'_r(\theta)|.
\]
The three-factor bound therefore proves \eqref{eq:C1}. Since
$S_{r+1}\le S_2=\pi^2/8-1<1/4$ and $1+\sqrt2<5/2$, one even has
$\eta_r<5/8<1$; its numerical upper value is about $0.5642$.

The first inequality identifies the zero set and sign; the second
identifies the stationary set and sign of the derivative. At a zero of
$\tau_r$, its derivative is nonzero, so \eqref{eq:C1} gives simplicity.
At a stationary point $\theta_0$, divide \eqref{eq:C1} by
$|\theta-\theta_0|$ and take the limit. The resulting bound
\[
 |\mathcal F''_{r,m}(\theta_0)-\tau''_r(\theta_0)|
 \le\eta_r|\tau''_r(\theta_0)|
\]
shows that the second derivative is nonzero there.
\end{proof}

\begin{corollary}[Baseline exactness classification]\label[corollary]{cor:phases}
At the mesh \eqref{eq:mesh}, the rule is exact on $\Poly_r$ if and only
if $\theta\in Z_r$. At every other phase its exact degree is $r-1$.
\end{corollary}
\begin{proof}
All lower monomial defects vanish. The remaining degree-$r$ defect is a
nonzero scalar times \eqref{eq:F}, and \cref{thm:shape} locates its zeros.
\end{proof}

\subsection{Quantitative stability under perturbation}

The relative derivative estimate excludes an instability that a uniform
function bound alone would miss: a very small, highly oscillatory
perturbation can introduce many zeros unless its derivative is also
controlled.

\begin{theorem}[Stable two-zero structure]\label{thm:stability}
Fix $0<\varepsilon<1/4$, and let $g$ be a real $C^1$ function on $\T$.
Suppose
\begin{align}
 \norm{g}_\infty&<(1-c_r)\sin(2\pi\varepsilon),\label{eq:stab0}\\
 \norm{g'}_\infty&<2\pi(1-\eta_r)\cos(2\pi\varepsilon).
 \label{eq:stab1}
\end{align}
Then $\mathcal F_{r,m}+g$ has exactly two zeros. Each lies in the
$\varepsilon$-neighborhood of a different point $z\in Z_r$, and its
circular distance from that point is at most
\begin{equation}\label{eq:displacement}
 \frac{\norm{g}_\infty}
 {2\pi(1-\eta_r)\cos(2\pi\varepsilon)}.
\end{equation}
Both zeros are simple.
\end{theorem}
\begin{proof}
Outside the two neighborhoods, $|\tau_r|\ge\sin(2\pi\varepsilon)$.
Equation~\eqref{eq:C0} and \eqref{eq:stab0} exclude a zero there. At the
two endpoints of each neighborhood, the unperturbed function has opposite
signs, and those signs survive the perturbation. Hence a zero exists in
each neighborhood.

Inside such a neighborhood, $|\tau'_r|\ge2\pi\cos(2\pi\varepsilon)$,
with constant sign. Equations~\eqref{eq:C1} and \eqref{eq:stab1} imply
that the perturbed derivative is nonzero with that sign throughout. This
gives uniqueness and simplicity. Apply the mean value theorem to the
unperturbed function between $z$ and its perturbed zero to obtain
\eqref{eq:displacement}.
\end{proof}

\section{An exact formula for the next defect}
\label{sec:next}

Set $p=r+1$. The first defect sees only odd frequencies; the next defect
sees odd frequencies and frequencies of valuation one. Treating these two
sets separately is the main structural step.

\begin{lemma}[The next Taylor coefficient at an odd alias]
\label[lemma]{lem:next-derivative}
If $N=Mn$ and $n$ is odd, then
\begin{equation}\label{eq:next-derivative}
 H^{(r+1)}(N)=\frac{(r+1)!}{N^r}
 \left(\frac rN H(mn/2)-2^{-r}H'(mn/2)\right).
\end{equation}
\end{lemma}
\begin{proof}
Write $z=N+w$. For each $j<r$,
\[
 \sinc\!\left(\frac{\pi(N+w)}{2^j}\right)
 =(-1)^{N/2^j}\frac wN
 \left(1-\frac wN+O(w^2)\right).
\]
The sine numerator has no quadratic term in $w$; the displayed linear
correction comes from the reciprocal denominator. Multiplying the first
$r$ factors, whose signs have product $-1$, gives
\[
 -\frac{w^r}{N^r}\left(1-\frac{rw}{N}+O(w^2)\right).
\]
The remaining product is
$H(z/2^r)=H(mn/2)+2^{-r}wH'(mn/2)+O(w^2)$.
Taking the coefficient of $w^{r+1}$ and multiplying by $(r+1)!$ proves
the formula. All expansions are justified by analyticity near $N$.
\end{proof}

\begin{theorem}[Odd aliases plus a refined first defect]
\label{thm:next-formula}
Define the signed, nonzero normalization
\begin{equation}\label{eq:K}
 K_{r,m}=\frac{2(r+1)!A_m}{(2\pi M)^{r+1}}.
\end{equation}
For every $\theta\in\T$,
\begin{equation}\label{eq:next-formula}
 \boxed{\displaystyle
 \frac{E_{r+1,M}(\theta)}{K_{r,m}}
 =\oddSum\left(\frac{r a_{m,n}}{n^{r+1}}
              -\frac{b_{m,n}}{n^r}\right)\tau_{r+1}(n\theta)
 -2^{-(r+1)}\mathcal F_{r+1,m}(2\theta).}
\end{equation}
The first term of the odd sum is
$(r+\beta_m)\tau_{r+1}(\theta)$.
\end{theorem}
\begin{proof}
Pairing the positive and negative odd frequencies in
\eqref{eq:poisson}, and using \cref{lem:next-derivative}, gives
\[
 \frac{2(r+1)!A_m}{(2\pi M)^{r+1}}
 \oddSum\left(\frac{r a_{m,n}}{n^{r+1}}
              -\frac{b_{m,n}}{n^r}\right)\tau_{r+1}(n\theta),
\]
because $M2^{-r}=m/2$. The entire even-frequency part is exactly
$E_{r+1,2M}(2\theta)$, by replacing $n$ with $2n$ in the Fourier
series. The mesh $2M=2^rm$ has first threshold $r+1$, so
\cref{prop:first} expresses that even part as
$-2^{-(r+1)}K_{r,m}\mathcal F_{r+1,m}(2\theta)$.
Finally $a_{m,1}=1$ and $b_{m,1}=-\beta_m$.
\end{proof}

\begin{corollary}[Parity cancellation of the even component]
\label[corollary]{cor:even-cancel}
At a selected phase $\theta\in Z_r$, the even component in
\eqref{eq:next-formula} vanishes identically when $r$ is even. When $r$
is odd, $2\theta$ is an integer and
\begin{equation}\label{eq:even-size}
 |\mathcal F_{r+1,m}(2\theta)|\le1+S_{r+3}.
\end{equation}
\end{corollary}
\begin{proof}
If $r$ is even, $\theta\in\{1/4,3/4\}$ and $r+1$ is odd. Every
harmonic of $\mathcal F_{r+1,m}$ is a sine, hence vanishes at
$2\theta\in\{1/2,3/2\}$. For odd $r$, use $2\theta\in\Z$ and
$|a_{m,n}|\le n^{-2}$ in the series.
\end{proof}

\section{Derivative domination with explicit constants}
\label{sec:derivative}

The next defect involves derivatives of $H$ at all odd half-integer
multiples, not only its values. Define
\begin{equation}\label{eq:T}
 P(t)=\sinc(\pi t)\sinc(\pi t/2),\qquad
 T(t)=\prod_{j=2}^{\infty}\sinc(\pi t/2^j),\qquad H=PT.
\end{equation}
The tail $T$ is the Fourier transform of a probability supported on
$[-1/4,1/4]$, so
\begin{equation}\label{eq:Tprime}
 |T'(t)|\le\frac\pi2\quad(t\in\R).
\end{equation}
Indeed, its random-variable representation is the part of $Y$ left after
removing its first two summands, and differentiation contributes at most
$2\pi$ times its support radius.

\begin{lemma}[Uniform derivative envelope along odd multiples]
\label[lemma]{lem:b-bound}
For a positive odd $m$, let
\begin{equation}\label{eq:D}
 D_m=\frac{m\pi}{4}\left(1+\frac1{|T(m/2)|}\right)+\frac23.
\end{equation}
This is finite, and for every odd $n\ge3$,
\begin{equation}\label{eq:b-bound}
 |b_{m,n}|\le\frac{D_m}{n^2}.
\end{equation}
\end{lemma}
\begin{proof}
Put $t_0=m/2$. All the factors at $t_0$ are nonzero. The two exact
factor ratios from \cref{lem:multiplier} give
$|P(nt_0)|=|P(t_0)|/n^2$. The same factorwise inequality gives
$|T(nt_0)|\le|T(t_0)|$.

At an odd half-integer,
\[
 \frac{P'(mn/2)}{P(mn/2)}
 =\frac\pi2\cot\!\left(\frac{\pi mn}{4}\right)-\frac4{mn}.
\]
The missing first cotangent is zero. The displayed cotangent has magnitude
one, hence the whole expression has magnitude at most
$\pi/2+4/(mn)$. Differentiating $H=PT$ and applying
\eqref{eq:Tprime} yields
\[
 \frac{|H'(mn/2)|}{|A_m|}
 \le\frac1{n^2}\left(\frac\pi2+\frac4{mn}
                         +\frac\pi{2|T(m/2)|}\right).
\]
Multiply by $m/2$ and use $2/n\le2/3$.
\end{proof}

\begin{lemma}[The dyadic constants have absolute bounds]
\label[lemma]{lem:dyadic-constants}
For $m=1$,
\begin{equation}\label{eq:dyadic-constants}
 A_1>0,\qquad \beta_1>1,\qquad D_1<\frac52.
\end{equation}
\end{lemma}
\begin{proof}
Every factor of $H(1/2)$ is positive. Each sinc factor is strictly
decreasing as a function of its positive argument in $(0,\pi/2]$.
The $j=0$ logarithmic derivative at $t=1/2$ is $-2$ and the remaining
logarithmic derivatives are negative. Therefore
$\beta_1=-\tfrac12H'(1/2)/H(1/2)>1$.

For the tail, $\sinc x\ge1-x^2/6$ for real $x$, and its arguments here
are positive and at most $\pi/8$. The elementary product inequality
$1-\prod_j(1-v_j)\le\sum_jv_j$ for $0\le v_j\le1$ gives
\[
 T(1/2)\ge1-\sum_{j=2}^{\infty}
 \frac{\pi^2}{24\,4^j}=1-\frac{\pi^2}{288}>\frac{24}{25}.
\]
For the final inequality, $\pi<22/7$ is more than sufficient. Hence
\[
 D_1<\frac{11}{14}\left(1+\frac{25}{24}\right)+\frac23
 =\frac{763}{336}<\frac52.
\]
\end{proof}

\begin{remark}[The odd part cannot be discarded]
The value $\beta_m$ need not be positive. Numerically,
$\beta_9\approx-10.8871$ and $\beta_{17}\approx-40.7636$.
Thus the leading coefficient $r+\beta_m$ can be small near particular
levels. The proof for $m=1$ must not be applied to all $m$ by silently
assuming $\beta_m>1$. The general criterion below uses its actual sign
and magnitude.
\end{remark}

\section{Sharp fixed-phase order at every dyadic mesh}
\label{sec:dyadic}

The following general error envelope is the bridge from Fourier analysis
to a nonvanishing theorem. Define
\begin{equation}\label{eq:R}
 R_{r,m}=rS_{r+3}+D_mS_{r+2}
 +\mathbf1_{\{r\text{ odd}\}}\,2^{-(r+1)}(1+S_{r+3}).
\end{equation}
The indicator retains the exact cancellation of the even-frequency part.

\begin{theorem}[A quantitative maximality criterion]
\label{thm:criterion}
For every $\theta\in Z_r$,
\begin{equation}\label{eq:remainder}
 \abs{\frac{E_{r+1,M}(\theta)}{K_{r,m}}
 -(r+\beta_m)\tau_{r+1}(\theta)}\le R_{r,m}.
\end{equation}
Consequently, if
\begin{equation}\label{eq:criterion}
 |r+\beta_m|>R_{r,m},
\end{equation}
then the exact degree of $Q_{M,\theta}$ is $r$ at each $\theta\in Z_r$
and $r-1$ at every other phase. In particular, its maximum over all phases
is $r$. At a selected phase,
\begin{align}
 |K_{r,m}|(|r+\beta_m|-R_{r,m})
 &\le |E_{r+1,M}(\theta)|\nonumber\\
 &\le |K_{r,m}|(|r+\beta_m|+R_{r,m}),\label{eq:two-sided}
\end{align}
and its sign is that of
$K_{r,m}(r+\beta_m)\tau_{r+1}(\theta)$.
\end{theorem}
\begin{proof}
At a selected phase, $|\tau_{r+1}(\theta)|=1$. Remove the $n=1$ term
from the odd sum in \cref{thm:next-formula}. Its remaining magnitude is
at most
\[
 \sum_{\substack{n\ge3\\n\text{ odd}}}
 \left(\frac{r}{n^{r+3}}+\frac{D_m}{n^{r+2}}\right)
 =rS_{r+3}+D_mS_{r+2},
\]
by \cref{lem:multiplier,lem:b-bound}. Add the even-component bound from
\cref{cor:even-cancel} to obtain \eqref{eq:remainder}. Under
\eqref{eq:criterion}, the leading term cannot be canceled. The two-sided
bound and sign follow, and \cref{cor:phases} completes the degree
classification.
\end{proof}

\begin{theorem}[Complete dyadic fixed-phase maximality]
\label{thm:dyadic}
For every $d\ge0$, put $M=2^d$ and $r=d+1$. Then
\begin{equation}\label{eq:degree-classification}
 \deg_{\rm exact}Q_{2^d,\theta}=
 \begin{cases}
 d+1,&\theta\in Z_{d+1},\\
 d,&\theta\notin Z_{d+1}.
 \end{cases}
\end{equation}
Here ``exact degree'' means the largest $s$ such that the rule is exact on
all of $\Poly_s$. More quantitatively, at either selected phase,
\begin{equation}\label{eq:quarter-margin}
 \abs{\frac{E_{r+1,2^{r-1}}(\theta)}
 {K_{r,1}(r+\beta_1)\tau_{r+1}(\theta)}-1}
 <\frac14.
\end{equation}
In particular, both selected phases fail at the very next degree.
\end{theorem}
\begin{proof}
By \cref{lem:dyadic-constants}, $r+\beta_1>r+1$ and $D_1<5/2$.
We may ignore the parity indicator for an upper bound. Therefore
\[
 \frac{R_{r,1}}{r+\beta_1}
 <\frac{rS_{r+3}+(5/2)S_{r+2}
              +2^{-(r+1)}(1+S_{r+3})}{r+1}.
\]
All terms are uniformly controlled by the first refinement level.
For $r\ge1$,
\[
 \frac{rS_{r+3}}{r+1}\le\frac12S_4,
 \qquad \frac{S_{r+2}}{r+1}\le\frac12S_3,
 \qquad \frac{2^{-(r+1)}(1+S_{r+3})}{r+1}
 \le\frac18(1+S_4).
\]
For the first inequality, use
$S_{r+3}\le3^{-(r-1)}S_4$ and
$r/(r+1)\le3^{r-1}/2$; the latter follows at $r=1$ by equality and
for $r\ge2$ from $r/(r+1)<1\le3^{r-1}/2$.
The other inequalities follow directly by monotonicity.

Equation~\eqref{eq:Sbound} gives
$S_3\le7/108$ and $S_4\le1/54$. Consequently
\begin{equation}\label{eq:rational-margin}
 \frac{R_{r,1}}{r+\beta_1}
 <\frac12S_4+\frac54S_3+\frac18(1+S_4)
 \le\frac{47}{216}<\frac14.
\end{equation}
Apply \cref{thm:criterion}; division by the signed leading term gives
\eqref{eq:quarter-margin}. The error cannot be zero, so the selected
phases have exact degree $r$. All other phases were already classified.
\end{proof}

\subsection{Why the result is stronger than phase completeness}

Knowing that the degree-$r$ error has exactly two zeros does not exclude
one of those zeros from also annihilating the degree-$(r+1)$ error.
The new ingredient is \eqref{eq:quarter-margin}: it separates the next
error from zero at \emph{both} selected phases, uniformly over all $d$.
The proof uses more than symmetry. Symmetry supplies selected zeros, but
only the derivative envelope and the refined-mesh decomposition control
the next error there.

\subsection{Signs and the two selected phases}

Since $A_1>0$ and $r+\beta_1>0$, the sign of the next dyadic defect is
precisely $\tau_{r+1}(\theta)$. At the pair of selected phases these
signs are opposite. Their magnitudes need not be equal when $r$ is odd:
the even-frequency component is then identical at the two phases and
shifts both errors by the same amount. When $r$ is even the errors are
exact negatives, because the even component vanishes and all remaining
frequencies are odd.

For example, the rule at $M=1$, $\theta=0$ has only the central nonzero
sample and gives $Q(x^2)=0$, whereas $I(x^2)=1/9$. At $\theta=1/2$,
the two samples at $\pm1/2$ give $Q(x^2)=1/4$, and the error is $5/36$.
This simple asymmetric pair is the first instance of the parity mechanism,
not an exceptional normalization artifact.

\section{Every fixed odd part has an effective finite frontier}
\label{sec:eventual}

\subsection{A cutoff from computable Fourier constants}

\begin{theorem}[Eventual maximality at every fixed odd part]
\label{thm:eventual}
For positive odd $m$, define the integer
\begin{equation}\label{eq:cutoff}
 r_*(m)=\max\left\{2,\ \left\lceil2|\beta_m|+2\right\rceil,
 \ 2+\left\lceil\log_3(1+D_m)\right\rceil\right\}.
\end{equation}
For every integer $r\ge r_*(m)$, the mesh $M=2^{r-1}m$ has exact
fixed-phase degree $r$ on $Z_r$ and $r-1$ off $Z_r$. In particular, any
failure of this classification for a fixed odd $m$ occurs at one of the
finitely many levels $1\le r<r_*(m)$.
\end{theorem}
\begin{proof}
For $s\ge3$, \eqref{eq:Sbound} gives $S_s\le2\cdot3^{-s}$.
Moreover, $S_{r+3}<1$, so
\begin{equation}\label{eq:Rcrude}
 R_{r,m}\le\frac{2r}{3^{r+3}}+
                \frac{2D_m}{3^{r+2}}+2^{-r}.
\end{equation}
Under \eqref{eq:cutoff}, $D_m<3^{r-2}$,
$r+\beta_m\ge r-|\beta_m|\ge r/2+1$, and $r\ge2$.
The first term in \eqref{eq:Rcrude} is at most $4/243$, since
$r/3^r$ decreases for integer $r\ge2$. The second is less than $2/81$
and the third at most $1/4$. Their sum is less than one, while the
leading coefficient is greater than one. Hence \eqref{eq:criterion}
holds, and \cref{thm:criterion} applies.
\end{proof}

The cutoff is intentionally sufficient, not optimal. It can be much larger
than the first level at which \eqref{eq:criterion} succeeds. A failed
sufficient test is not evidence that the underlying error is zero.

\subsection{A fully elementary cutoff, without special-function constants}

To remove even the need to evaluate $\beta_m$ and $D_m$, we bound them
using dyadic separation from the sine zeros.

\begin{lemma}[Elementary bounds on the odd-part constants]
\label[lemma]{lem:elementary}
For odd $m\ge1$ let
\begin{equation}\label{eq:J}
 J=\max\{2,1+\lceil\log_2m\rceil\}.
\end{equation}
Then
\begin{equation}\label{eq:elementary-constants}
 |\beta_m|<J(m+1)+1,\qquad
 D_m<m+(2m)^{J-1}+1\le(2m)^J.
\end{equation}
\end{lemma}
\begin{proof}
Set $t=m/2$ and $c_j=\pi/2^j$. Since $2^J\ge2m$ and $\pi<4$,
$|c_Jt|<1$. At an odd half-integer the angle
$c_jt=\pi m/2^{j+1}$ has distance at least $\pi/2^{j+1}$ from the
nearest integer multiple of $\pi$. Since $|\cot x|\le1/\delta$
when that distance is $\delta\in(0,\pi/2]$,
\[
 |c_j\cot(c_jt)|\le2,\qquad
 |c_j\cot(c_jt)-1/t|\le2+2/m.
\]
Thus the first $J$ contributions to
$|tH'(t)/H(t)|$ total at most $J(m+1)$.

For $|x|\le1$, the bounds
$|\sin x|\ge5|x|/6$ and
$|x\cos x-\sin x|\le|x|^3/3$ give
\begin{equation}\label{eq:cot-tail}
 |x\cot x-1|\le\frac25x^2.
\end{equation}
The numerator bound follows by integrating $-x\sin x$ from zero.
Consequently the logarithmic-derivative tail contributes at most
\[
 \frac25t^2\sum_{j=J}^{\infty}c_j^2
 =\frac{8\pi^2t^2}{15}\,4^{-J}
 =\frac{2\pi^2m^2}{15}\,4^{-J}
 <\frac8{15}<1.
\]
This proves the bound for $|\beta_m|$.

For the bound on $D_m$, every factor of $T(t)$ with $2\le j<J$ has
magnitude at least $2/(\pi m)>1/(2m)$. Indeed, its sine numerator has
magnitude at least $\sin(\pi/2^{j+1})\ge2^{-j}$, whereas its argument
has magnitude $\pi m/2^{j+1}$. For $j\ge J$ all factors are positive.
The same product estimate used in \cref{lem:dyadic-constants} gives
\[
 \prod_{j=J}^{\infty}\sinc(c_jt)
 \ge1-\frac{2\pi^2t^2}{9}\,4^{-J}
 =1-\frac{\pi^2m^2}{18}\,4^{-J}>\frac79.
\]
It follows that
\[
 |T(m/2)|>\frac79(2m)^{-(J-2)}.
\]
Using $\pi/4<1$, $9/7<2$, and $2/3<1$ in \eqref{eq:D} gives
\[
 D_m<m\bigl(1+2(2m)^{J-2}\bigr)+1
 =m+(2m)^{J-1}+1.
\]
The last expression is at most $(2m)^J$: since $J\ge2$,
$(2m)^{J-1}(2m-1)\ge2m(2m-1)\ge m+1$ for $m\ge1$.
\end{proof}

\begin{theorem}[Explicit finite-level reduction]
\label{thm:elementary}
For every odd $m\ge1$, with $J$ as in \eqref{eq:J}, all refinement levels
\begin{equation}\label{eq:elementary-cutoff}
 r\ge 2J(m+1)+4
\end{equation}
have the sharp fixed-phase classification of \cref{thm:criterion}.
Equivalently, at fixed odd part $m$, it suffices to settle the finitely many
levels
\begin{equation}\label{eq:finite-levels}
 1\le r\le 2J(m+1)+3
\end{equation}
to complete the all-level classification for that odd part.
\end{theorem}
\begin{proof}
The $\beta_m$ bound in \cref{lem:elementary} makes
$2|\beta_m|+2<2J(m+1)+4$. Also
\[
 1+D_m<1+(2m)^J\le2(2m)^J\le2\cdot3^{Jm}<3^{Jm+1},
\]
where $2m\le3^m$ for integers $m\ge1$. Thus
$2+\lceil\log_3(1+D_m)\rceil\le Jm+3$, which is smaller than the
right side of \eqref{eq:elementary-cutoff}. The cutoff in
\cref{thm:eventual} is therefore satisfied.
\end{proof}

This reduction is finite for each prescribed $m$, not uniformly finite
over all odd integers. The possible number of unresolved levels still
grows with $m$. In particular, \cref{thm:elementary} is not a proof that
there are only finitely many exceptional integer meshes altogether.

\subsection{Effective evaluation and rigorous numerical certificates}

The constants in \eqref{eq:cutoff} are computable with certified errors.
For $t>0$, choose $J$ so that $\pi t2^{-J}\le1$. The product tail has
\begin{equation}\label{eq:product-tail}
 0\le1-\prod_{j=J}^{\infty}\sinc(\pi t/2^j)
 \le\frac{2\pi^2t^2}{9}\,4^{-J}.
\end{equation}
The logarithmic-derivative tail obeys, by \eqref{eq:cot-tail},
\begin{equation}\label{eq:log-tail}
 \abs{\sum_{j=J}^{\infty}
 \left(\frac\pi{2^j}\cot\!\left(\frac{\pi t}{2^j}\right)-\frac1t\right)}
 \le\frac{8\pi^2t}{15}\,4^{-J}.
\end{equation}
At an odd half-integer, the finite initial product is nonzero. These tail
bounds therefore provide a positive lower bound on its absolute value
once the cutoff is sufficiently large. Interval bounds on the finite
sines, cosines, and $\pi$ then enclose $\beta_m$ and $D_m$.

There is no algorithmic need to decide whether a transcendental expression
is an integer before choosing a cutoff: rational upper enclosures
$B>|\beta_m|$ and $D>D_m$ give a conservative integer cutoff by rounding
up. The completely elementary cutoff is available without these
computations. The supplied numerical implementation illustrates the
sharper constants, but does not implement directed-rounding intervals;
its output is explicitly labeled diagnostic.

\section{The two-exponential structure of the next error}
\label{sec:asymptotic}

The exact formula contains more information than eventual nonvanishing.
It distinguishes the first even alias from the first nontrivial odd
alias, and shows why their relative importance depends on parity.

For $s>1$ write
\[
 U_s=\sum_{\substack{n\ge5\\n\text{ odd}}}n^{-s}.
\]
As in \eqref{eq:Sbound},
\begin{equation}\label{eq:Ubound}
 U_s\le5^{-s}\left(1+\frac5{2(s-1)}\right).
\end{equation}

\begin{theorem}[Explicit three-term expansion at the selected phases]
\label{thm:expansion}
For $\theta\in Z_r$ and fixed odd $m$,
\begin{align}
 \frac{E_{r+1,M}(\theta)}{K_{r,m}}
 ={}&(r+\beta_m)\tau_{r+1}(\theta)
       -2^{-(r+1)}\tau_{r+1}(2\theta)\nonumber\\
 &+3^{-r}\left(\frac r3a_{m,3}-b_{m,3}\right)
              \tau_{r+1}(3\theta)+\mathcal R_{r,m}(\theta),
 \label{eq:expansion}
\end{align}
where
\begin{equation}\label{eq:expansion-remainder}
 |\mathcal R_{r,m}(\theta)|
 \le rU_{r+3}+D_mU_{r+2}
 +\mathbf1_{\{r\text{ odd}\}}2^{-(r+1)}S_{r+3}.
\end{equation}
Thus, as $r\to\infty$ with $m$ fixed, the remainder is
$O_m(r5^{-r}+6^{-r})$, uniformly over the two selected phases.
When $r$ is even, every even-alias term, including its remainder, is
exactly zero.
\end{theorem}
\begin{proof}
In the odd part of \eqref{eq:next-formula}, retain $n=1$ and $n=3$.
Their contributions are exactly the first and third terms of
\eqref{eq:expansion}. The remaining odd terms are bounded by
$rU_{r+3}+D_mU_{r+2}$. In the refined term retain only its $n=1$
harmonic. Its remaining terms have magnitude at most
$2^{-(r+1)}S_{r+3}$, unless $r$ is even, when all its harmonics vanish
by \cref{cor:even-cancel}. Finally use \eqref{eq:Ubound} and
\eqref{eq:Sbound}.
\end{proof}

\begin{corollary}[Signed leading asymptotic]\label[corollary]{cor:leading}
For fixed odd $m$, uniformly over $\theta\in Z_r$,
\begin{equation}\label{eq:leading}
 E_{r+1,2^{r-1}m}(\theta)
 =K_{r,m}(r+\beta_m)\tau_{r+1}(\theta)
 \left(1+O_m\!\left(3^{-r}+\frac{2^{-r}}r\right)\right).
\end{equation}
On the even-$r$ subsequence the term $2^{-r}/r$ may be omitted.
\end{corollary}
\begin{proof}
For fixed $m$, $r+\beta_m\sim r$. Divide the error estimate in
\cref{thm:criterion} by $|r+\beta_m|$, and use
\eqref{eq:Sbound}. The even-$r$ improvement follows from the exact
cancellation.
\end{proof}

The two small parameters are not interchangeable. For odd $r$, the
$2^{-r}$ term occurs before the $r3^{-r}$ odd correction. For even $r$,
the whole $2^{-r}$ family disappears. This is a simple instance of an
exponential asymptotic scale selected by arithmetic parity, directly
relevant to the repository's broader interest in transseries. No divergent
series, summability prescription, or beyond-all-orders assertion is needed
here: the finite truncation has an explicit remainder bound.

\subsection{Exact sum and difference of the two errors}

Let $\theta_0$ and $\theta_0+1/2$ be the selected phases. Every odd
harmonic changes sign under this shift, whereas the refined component is
unchanged. Hence the identity
\begin{equation}\label{eq:pair-sum}
 E_{r+1,M}(\theta_0)+E_{r+1,M}(\theta_0+1/2)
 =2E_{r+1,2M}(2\theta_0)
\end{equation}
is exact. Their difference isolates twice the odd part of
\eqref{eq:next-formula}. In particular, if $r$ is even their sum is zero.
If $r$ is odd, the sum has the sign and size of a first defect on the
refined mesh. This supplies a direct check on both the factor $2^{-(r+1)}$
and its minus sign in the main formula.

\section{Polynomial synthesis and refined phase filters}
\label{sec:synthesis}

\subsection{Moment deconvolution preserves the sharp degree}

Define the polynomial averaging operator
\begin{equation}\label{eq:B}
 (\mathcal BP)(t)=\int_\R P(t-x)u(x)\dd x.
\end{equation}
In the monomial basis it is triangular with diagonal one, so it has an
inverse $\mathcal D$ on the polynomial algebra. Both operators preserve
degree and leading coefficient. The polynomials $A_n=\mathcal D(t^n)$
are monic and have the formal generating identity
\[
 \sum_{n\ge0}A_n(t)\frac{z^n}{n!}
 =\frac{e^{tz}}{\int_\R e^{zx}u(x)\dd x}.
\]
Evenness of $u$ identifies its moment-generating function at $z$ and
$-z$. Only a formal power series, or a neighborhood of zero, is intended;
the reciprocal denominator is not asserted to be entire.

The phased synthesis operator is
\begin{equation}\label{eq:synthesis}
 (\mathcal S_{M,\theta}P)(t)
 =\frac1M\sum_{k\in\Z}
 (\mathcal DP)\!\left(t-\frac{k+\theta}{M}\right)
 u\!\left(\frac{k+\theta}{M}\right).
\end{equation}
This is the deconvolution framework already present in the repository
\cite{repo-lean,repo-phase}.

\begin{theorem}[Sharp synthesis order and its leading failure]
\label{thm:synthesis}
The identity $\mathcal S_{M,\theta}P=P$ for every $P\in\Poly_s$ is
equivalent to exactness of $Q_{M,\theta}$ on $\Poly_s$.
If the quadrature is exact through degree $r$ and $p=r+1$, then
\begin{equation}\label{eq:synthesis-error}
 (\mathcal S_{M,\theta}t^p)(t)-t^p=(-1)^pE_{p,M}(\theta),
\end{equation}
independently of $t$. More generally the error for a degree-at-most-$p$
polynomial is its leading degree-$p$ coefficient times the right side.
Consequently all the sharp-order conclusions of
\cref{thm:dyadic,thm:eventual,thm:elementary} apply to phased synthesis.
\end{theorem}
\begin{proof}
Quadrature exactness applies to the polynomial
$x\mapsto(\mathcal DP)(t-x)$, whose integral is
$(\mathcal B\mathcal DP)(t)=P(t)$. This proves one direction. Conversely,
as $P$ ranges over $\Poly_s$, the polynomials
$x\mapsto(\mathcal DP)(-x)$ range over all of $\Poly_s$, since
$\mathcal D$ and reflection are invertible. Evaluating synthesis at $t=0$
therefore implies quadrature exactness.

For $P=t^p$, the polynomial $A_p(t-x)$ has leading coefficient $(-1)^p$
as a polynomial in $x$. Every lower-degree quadrature defect vanishes,
so only this coefficient multiplies $E_{p,M}$. The general statement
follows by linearity and preservation of leading coefficients.
\end{proof}

\subsection{Consequences for phase averages}

For any integer $L\ge1$ and offset $\alpha$,
\begin{equation}\label{eq:refinement}
 \frac1L\sum_{j=0}^{L-1}Q_{M,\theta+\alpha+j/L}(P)
 =Q_{LM,L(\theta+\alpha)}(P).
\end{equation}
This is simply the bijection $(k,j)\mapsto Lk+j$ between the lattice
indices, and holds beyond polynomials whenever the sums exist. The earlier
phase-filter manuscript identifies such uniform refinements as the sharp
minimum-size filters for uniform translated exactness \cite{repo-phase}.

\begin{corollary}[Fixed-phase order of a dyadic refinement]
\label[corollary]{cor:refinement}
If $LM=2^d$, the phase average in \eqref{eq:refinement} has exact degree
$d+1$ precisely at the translating phases satisfying
\[
 L(\theta+\alpha)\in Z_{d+1}\pmod1.
\]
It has exact degree $d$ at every other translating phase. There are exactly
$2L$ maximizing phases modulo one, and no translating phase gives degree
$d+2$. If $LM=2^{r-1}m$ instead, the same conclusion with $d=r-1$
holds whenever \eqref{eq:criterion} or either sufficient cutoff applies.
\end{corollary}
\begin{proof}
Apply the exact refinement identity and the corresponding sharp-order
theorem. Multiplication by $L$ on $\T$ has exactly $L$ preimages of each
of the two selected phases.
\end{proof}

This corollary adds a fixed-phase ceiling to an existing uniform-filter
classification. It neither changes that classification nor asserts
optimality among arbitrary unequal-weight filters at a single phase.

\section{Independent exact and numerical verification}
\label{sec:verification}

The package includes a standalone Python program,
\path{reproducibility/verify.py}. It imports no repository code. Its exact
component uses \code{fractions.Fraction}; its numerical component uses
\code{mpmath} with 60 decimal digits. The two kinds of evidence are recorded
separately and are not conflated with the proofs above.

\subsection{Exact dyadic evaluator and moment computation}

For the bounded Fabius function $F$ on $[0,1]$, the normalization is
$u(x)=F(1-|x|)$ for $|x|\le1$. The classical dyadic recurrences from
\cite{arias2017}, also used by the earlier repository manuscript, give an
exact evaluator. Put
\begin{equation}\label{eq:dn}
 d_n=n!2^{n(n-1)/2}F(2^{-n}),\qquad d_0=1.
\end{equation}
Then
\begin{equation}\label{eq:drec}
 (n+1)(2^n-1)d_n=\sum_{k=0}^{n-1}\binom{n+1}{k}d_k\qquad(n\ge1).
\end{equation}
For a dyadic $x\in(0,1)$, select $n$ so that
$2^{-n}\le x<2^{1-n}$ and put $y=x-2^{-n}$. The identity
\begin{equation}\label{eq:Frec}
 F(x)=-F(y)+\sum_{k=0}^{n}
 \frac{2^{k(k+1)/2}}{k!}F(2^{k-n})y^k
\end{equation}
removes the leading nonzero binary digit. Values at powers of two are
computed directly from \eqref{eq:dn}, preventing a circular recursive
call. The endpoint values are $F(0)=0$ and $F(1)=1$.

Moments are computed independently from the distributional identity
$Y\overset d=(V+Y')/2$, with $V$ uniform on $[-1,1]$ and independent
of a copy $Y'$. With $\mu_n=I(x^n)$, odd moments vanish and
\begin{equation}\label{eq:moments}
 (2^n-1)\mu_n=
 \sum_{\substack{0\le k<n\\k\text{ even}}}
 \binom nk\frac{\mu_k}{n-k+1}
 \qquad(n>0\text{ even}),\qquad\mu_0=1.
\end{equation}
This follows by expanding $(V+Y')^n$ and isolating the last term.
The evaluator and moment recurrence therefore approach the quadrature
identity from different finite computations.

\begin{table}[htbp]
\centering
\caption{Exact next-degree failures at both maximizing phases. The general
theorem, not these finitely many rows, proves all-level maximality.}
\label{tab:exact}
\renewcommand{\arraystretch}{1.4}
\begin{tabular}{@{}ccccc@{}}
\toprule
$M$ & $r$ & Selected phases & First next defect & Second next defect\\
\midrule
$1$ & $1$ & $0,\ 1/2$ & $-1/9$ & $5/36$\\
$2$ & $2$ & $1/4,\ 3/4$ & $17/1536$ & $-17/1536$\\
$4$ & $3$ & $0,\ 1/2$ & $97/345600$ & $-799/2764800$\\
$8$ & $4$ & $1/4,\ 3/4$ & $-9919/4529848320$ & $9919/4529848320$\\
\bottomrule
\end{tabular}
\end{table}

The exact run performs 978 assertions. These include 257 reflection checks
for $F$; lower-degree exactness and first-threshold phase classification on
a 16-point phase grid at each mesh $1,2,4,8,16,32,64$; next-degree
nonvanishing at both selected phases; the four comparison values in
\cref{tab:exact}; and 126 exact lattice-refinement identities. Every
assertion passed. The complete rational outputs, including the larger
fractions through $M=64$, are in \path{validation/exact_results.json}.

These exact computations verify finite examples, not infinite-dimensional
analytic statements. In particular, testing a finite phase grid does not
prove that other phases are absent; \cref{thm:shape} supplies that proof.

\subsection{Product and derivative diagnostics}

The numerical program computes $H$ and $H'$ from 112 product factors and
the corresponding logarithmic-derivative sum. For small arguments it uses
Taylor expansions to avoid subtracting nearly equal numbers. It tests both
bounds in \eqref{eq:multiplier} and the derivative bound
\eqref{eq:b-bound} for
\[
 m\in\{1,3,5,7,9,11,13,15,17,31\},\qquad
 n\in\{3,5,\ldots,63\}.
\]
These give 930 checks. Two additional checks verify the numerical values
of the inequalities for $\beta_1$ and $D_1$.

For each of the fourteen selected-phase cases through $M=64$, the program
compares the exact rational next defect with \eqref{eq:next-formula},
truncated at odd index 201, and also checks the relative $1/4$ bound.
That gives another 28 checks, for 960 numerical assertions in total.
All passed. The largest discrepancy in the \emph{normalized} next-defect
formula was approximately $2.7651\cdot10^{-15}$, and the largest observed
relative deviation from its leading term was approximately $0.132868$.
The latter occurs at the smallest mesh in the test.

These values describe this particular numerical run. They are not
rigorous decimal enclosures. The checks compare series discrepancies to
an analytic tail envelope, but the arithmetic itself does not use directed
rounding. The proofs of the envelope and of the $1/4$ margin are separate.
The detailed output is in \path{validation/numerical_results.json}.

\begin{table}[htbp]
\centering
\caption{Odd-part constants and conservative cutoffs. Decimal constants and
the $r_*$ column are numerical diagnostics; the last column follows from the
proved elementary formula. For $m=1$, the sharper dyadic theorem already
covers every $r\ge1$.}
\label{tab:constants}
\begin{tabular}{@{}rrrrr@{}}
\toprule
$m$ & $\beta_m$ (approx.) & $D_m$ (approx.) & $r_*(m)$ (approx.)
 & $2J(m+1)+4$\\
\midrule
1 & 1.283712 & 2.264960 & 5 & 12\\
3 & 5.025370 & 6.270594 & 13 & 28\\
5 & 0.338832 & 14.993681 & 5 & 52\\
7 & 16.076234 & 67.714799 & 35 & 68\\
9 & $-10.887077$ & 148.173749 & 24 & 104\\
13 & 0.522487 & 389.713427 & 8 & 144\\
17 & $-40.763598$ & 7839.404861 & 84 & 220\\
31 & 121.429034 & 1302794.266702 & 245 & 388\\
\bottomrule
\end{tabular}
\end{table}

The sufficient inequality itself often succeeds much earlier than either
cutoff. In the numerical run, it succeeded at every tested level
$r=1,\ldots,20$ for $m=1,3,5,7,9,11$. For $m=13,15,17$ it succeeded
at $r=3,\ldots,20$, and for $m=31$ at $r=7,\ldots,20$. These are
observations about a sufficient test, not certified complete
classifications of those odd parts. The elementary all-large-level theorem
requires no reliance on this table.

\subsection{Reproduction and validation boundaries}

The exact checks need only Python's standard library. Run
\begin{quote}\small\ttfamily
python reproducibility/verify.py --exact-only
\end{quote}
The numerical checks additionally need \code{mpmath}:
\begin{quote}\small\ttfamily
python reproducibility/verify.py --numeric-only
\end{quote}
Running without either switch performs both sets. The script raises an
exception at a failed assertion; it does not silently accept failed tests.
The JSON files preserve exact rational strings separately from rounded
floating-point displays.

No Lean build was performed. No numerical claim here is described as a
formal theorem merely because a Python assertion passed. Conversely, the
analytic theorems do not assume a successful test run: their proofs are
contained in the preceding sections.

\section{Formalization plan and a proof-dependency boundary}
\label{sec:formalization}

\subsection{What can be reused}

The inspected repository already has the Poisson-defect framework and
sufficient parity-phase exactness \cite{repo-defect,repo-lean}. Relevant declarations include
\begin{quote}\small\ttfamily
summable\_fourier\_monomialRvachevSchwartz\par\medskip
tsum\_shifted\_monomial\_sub\_integral\par\medskip
tsum\_shifted\_monomial\_sub\_integral\_odd\par\medskip
tsum\_shifted\_polynomial\_eq\_integral\_superconvergent.
\end{quote}
These names are dependency candidates, not claims that the new results
already exist in Lean. The polynomial synthesis layer can reuse the
existing deconvolution framework. The classical analytic normalization and
zero-order facts should be matched carefully before any new theorem is
stated; shifting the Fourier product's starting index changes the support
and every mesh factor.

\subsection{A staged formal development}

A useful first target is \cref{thm:dyadic}, without first formalizing the
more general odd-part cutoff. The minimal new chain is: obtain the local
Taylor coefficient at a zero; split the next defect into odd and even
frequencies; prove the derivative envelope with $m=1$; establish the
rational bounds $\beta_1>1$ and $D_1<5/2$; and combine the summable tails
with \eqref{eq:rational-margin}.

The proof does not need a full zeta-function library. The bounds
$S_3\le7/108$ and $S_4\le1/54$, together with the monotonic comparison
used in the proof, suffice for dyadic nonvanishing. Similarly,
$\eta_r<5/8$ suffices for the stationary-point theorem; the decimal
constant is not needed.

A second stage can generalize the derivative envelope and prove
\cref{thm:criterion}. A third stage can add the integer-valued elementary
cutoff. In that final stage it may be cleaner to define $J$ as the least
integer at least two satisfying $2^J\ge2m$, rather than introduce real
logarithms and ceilings. This is equivalent to \eqref{eq:J}, but exposes
exactly the inequality used by the proof.

The asymptotic expansion should come only after the exact formula and
finite tail inequalities. Its useful formal core is not a notation for
big-$O$, but the explicit bound \eqref{eq:expansion-remainder}. Finally,
the $C^1$ perturbation theorem can be proved separately as a general result
about a harmonic plus a controlled perturbation.

\subsection{Suggested endpoint statements}

For dyadic meshes, the natural endpoint is the pair of statements
\[
 \forall\theta\in Z_{d+1},\quad E_{d+2,2^d}(\theta)\ne0,
\]
and
\[
 \forall\theta\in\T,\quad
 \bigl[Q_{2^d,\theta}\text{ exact on }\Poly_{d+1}\bigr]
 \Longleftrightarrow\theta\in Z_{d+1}.
\]
The first is the new maximality content; the second is the already known
phase-completeness boundary. Their combination yields
\eqref{eq:degree-classification}. Working on the circle, or explicitly
reducing a real phase modulo one, is essential.

For the general theorem, a formal statement should retain $m>0$, oddness,
$r\ge1$, and the decomposition $M=2^{r-1}m$. The sign of $A_m$ and
$\beta_m$ must not be discarded. Using an absolute normalization for the
next error would require an extra sign factor and makes the exact identity
less transparent; \eqref{eq:K} deliberately keeps the signed normalization.

\section{Further research questions and topics}
\label{sec:research}

The results give several concrete directions with explicit success
criteria. The following questions are proposed, not solved by the present
manuscript.

\subsection{Complete the all-integer-mesh maximality classification}

Does $E_{r+1,2^{r-1}m}(\theta)$ ever vanish at a phase $\theta\in Z_r$
for any positive odd $m$ and any $r\ge1$? The dyadic case is now excluded,
and \cref{thm:elementary} excludes all sufficiently large levels for each
fixed odd part. What remains is not a single finite computation, because
$m$ is unbounded. A complete theorem would require arithmetic information
that is uniform in the odd part, or an exact counterexample.

A practical first milestone is an interval-certified classification for
all odd $m\le B$, checking only the finitely many levels from
\eqref{eq:finite-levels}. This would turn the new theorem into a rigorous
finite-frontier search rather than an arbitrary grid experiment.

\subsection{Control near-resonance of the leading coefficient}

The obstruction to a one-line extension of the dyadic proof is the
possible smallness of $r+\beta_m$. Determine quantitative lower bounds on
its distance from zero, or classify its near-integer values. The exact
identity suggests that near-resonance may force the $n=3$ harmonic or the
even refined term to become decisive. A proof that these terms cannot
simultaneously cancel would go beyond first-harmonic domination.

\subsection{Replace the conservative derivative envelope}

The term $1/|T(m/2)|$ in $D_m$ can be very large. It results from bounding
$|T'(mn/2)|$ by a global support estimate, which ignores the strong
arithmetic relation between $mn/2$ and $m/2$. A relative derivative
inequality along odd multiples could substantially improve the cutoff.
The target is an estimate for $b_{m,n}$ that retains cancellation in the
logarithmic derivative rather than separately bounding $P'T$ and $PT'$.

\subsection{Find the optimal dyadic relative margin}

The constant $1/4$ in \eqref{eq:quarter-margin} is convenient, not sharp.
The exact check at $M=1$, $\theta=0$ gives a relative deviation around
$0.132868$, while the proved upper bound is $47/216\approx0.217593$.
Is the supremum over all dyadic levels and both selected phases attained
at that smallest mesh? Proving such a statement would require a sharper
comparison between consecutive levels, not merely separate tail bounds.

\subsection{Higher derivatives and the geometry of later defects}

The three-factor estimate controls one phase derivative. Additional fixed
factors give further powers of odd-multiplier decay, but their residue
constants grow. Determine which higher derivative orders admit a uniform
harmonic-domination theorem independent of $m$ and $r$. Separately,
classify the zero sets and stationary points of the next and later
monomial defects. Their even aliases create a geometry absent from the
first defect.

\subsection{All-order alias expansions with effective remainders}

Retaining any finite set of odd frequencies in \eqref{eq:next-formula}
gives additional exponential scales. For higher monomial degrees, several
valuation layers survive. Derive a systematic expansion indexed jointly
by frequency and its two-adic valuation, with explicit coefficients from
local Taylor jets of $H$ and uniform truncation bounds. The present
$2^{-r}$/$3^{-r}$ split is the first nontrivial case of that program.

\subsection{Generalized geometric convolution kernels}

Replace the dyadic uniform series by other geometrically weighted
convolutions. Which arithmetic properties of their Fourier zeros permit
complete phase classification and fixed-phase ceilings? Integer bases
larger than two have more residue classes and more candidate phases.
One should distinguish what follows from zero multiplicities alone from
what needs special sine-ratio estimates like \eqref{eq:multiplier}.

\subsection{Certified filters robust to phase and weight errors}

The perturbation theorem controls the normalized first defect when a
concrete implementation can bound its $C^1$ perturbation. Develop explicit
bounds translating errors in sampled values, phase locations, and filter
weights into those norms. Combining such bounds with the next-degree
nonvanishing margin would provide a certified region in which both the
number of superconvergent phases and the fixed-phase ceiling survive.
This is a stronger numerical-design question than merely observing stable
plots.

\section{Conclusion}

The repository's first-superconvergence phase classification leaves a
logically separate question: whether a special phase might also cancel
the next monomial error. The exact odd/even decomposition answers that
question at every dyadic mesh, with a uniform relative nonvanishing
margin. The same mechanism produces effective eventual maximality for
every fixed odd part, including a completely elementary finite-level
reduction.

The proof also explains the geometry and scale separation of the errors.
A third sinc factor controls the derivative of the normalized first
defect, while the next defect divides into a same-mesh odd series and a
refined first defect. These two observations supply a foundation for
further exact classification, certified computation, and formalization,
without treating numerical checks as proofs or silently extending a
dyadic statement to all integer meshes.

\appendix
\section{A compact dependency and novelty ledger}
\label{app:ledger}

\begin{longtable}{@{}p{0.30\textwidth}p{0.63\textwidth}@{}}
\toprule
Result or ingredient & Status and role\\
\midrule
\endhead
Fourier product and zero multiplicities & Classical kernel structure;
\cref{sec:fourier} fixes normalization and supplies a direct proof of the
part used here.\\[5pt]
Two-factor multiplier bound and first-phase completeness & Present in
\cite{repo-phase}; re-established as the baseline, not claimed new.\\[5pt]
Three-factor multiplier bound and complete stationary-point geometry &
Developed here in \cref{lem:multiplier,thm:shape}; gives a relative $C^1$
control, rather than only the first zero set.\\[5pt]
Exact next-defect decomposition & Developed in
\cref{thm:next-formula}; keeps the derivative term and the entire
refined even-frequency component.\\[5pt]
Sharp dyadic fixed-phase degree & Main completion in
\cref{thm:dyadic}; covers every level, both selected phases, and gives a
uniform signed relative margin.\\[5pt]
General odd-part maximality & Sufficient criterion and eventual theorem in
\cref{thm:criterion,thm:eventual}; not an unconditional all-level theorem
for every odd $m$.\\[5pt]
Elementary finite-level reduction & \Cref{thm:elementary} gives the explicit
$O(m\log(m+1))$ cutoff without evaluating a special-function constant.\\[5pt]
Parity-dependent exponential expansion & \Cref{thm:expansion} gives exact
retained terms and a finite-error inequality.\\[5pt]
Synthesis and uniform-grid refinement & Existing frameworks;
\cref{thm:synthesis,cor:refinement} transfer the new sharp ceilings.\\[5pt]
Formal verification & No new Lean theorem or Lean build is claimed.
The proof boundary and staged plan are given in \cref{sec:formalization}.\\
\bottomrule
\end{longtable}

\section{Normalization and edge-case checklist}

The first Fourier factor is $\sinc(\pi t)$, not
$\sinc(\pi t/2)$. The support is $[-1,1]$, the integral is one, and
$u(0)=1$. The latter follows by convolving the first uniform summand,
whose density equals one on $[-1/2,1/2]$, with the remaining law supported
on that interval. It is also part of the classical normalization.

The case $M=1$ is included: $r=1$ and the base exact degree is zero.
The phase circle identifies $0$ with $1$; no theorem excludes integer
translates of a selected representative. The odd part is positive; negative
frequencies are handled by evenness of $H$, not by treating $m$ as signed.
The factor $A_m$ is nonzero but may be negative. Every next-error formula
retains its sign. The integer cutoff uses $r=d+1$, so a cutoff on $r$
becomes a cutoff one smaller on $d$.

Finally, exactness through degree $s$ is a simultaneous statement about
all polynomials of degree at most $s$. A single higher monomial can
sometimes be integrated accidentally even when lower ones fail. Such an
isolated equality does not increase the exact degree used in this article.

\begin{thebibliography}{9}

\bibitem{arias1982}
J.~Arias de Reyna,
\emph{An infinitely differentiable function with compact support:
Definition and properties}, English translation (2017) of
\emph{Definici\'on y estudio de una funci\'on indefinidamente diferenciable
de soporte compacto}, Rev.~Real Acad.~Ciencias \textbf{76} (1982), 21--38.
\url{https://arxiv.org/abs/1702.05442}.

\bibitem{arias2017}
J.~Arias de Reyna,
\emph{Arithmetic of the Fabius function}, arXiv:1702.06487, version 3 (2017).
\url{https://arxiv.org/abs/1702.06487}.

\bibitem{repo-phase}
ProveIt repository,
\emph{Complete Superconvergence Phases and Optimal Phase Filters for
Rvachev Quadrature}, research manuscript, 28 September 2026.
Inspected at commit
\texttt{23dd71d2d3d5e85d2b97a8cc45f55e735d69a4ae}; source:
\begin{small}
\url{https://github.com/VladimirReshetnikov/ProveIt/blob/23dd71d2d3d5e85d2b97a8cc45f55e735d69a4ae/Analysis/FabiusFunction/docs/semi-formalized-research-frontiers/drafts/inverse-and-sampling/comb-interpolation/Optimal_Phase_Filters/article.tex}.
\end{small}
The introduction and research discussion distinguish first-phase
completeness from the unresolved fixed-phase next-degree maximality
question. This reference is an unrefereed repository manuscript.

\bibitem{repo-lean}
ProveIt repository,
\code{RvachevSuperconvergentSynthesis.lean}, inspected at the same commit.
\begin{small}
\url{https://github.com/VladimirReshetnikov/ProveIt/blob/23dd71d2d3d5e85d2b97a8cc45f55e735d69a4ae/Analysis/FabiusFunction/Lean/FabiusFunction/RvachevSuperconvergentSynthesis.lean}.
\end{small}

\bibitem{repo-defect}
ProveIt repository,
\code{CombDefectSeries.lean}, dependency named by the inspected
superconvergence module and the phase-filter manuscript.
\begin{small}
\url{https://github.com/VladimirReshetnikov/ProveIt/blob/23dd71d2d3d5e85d2b97a8cc45f55e735d69a4ae/Analysis/FabiusFunction/Lean/FabiusFunction/CombDefectSeries.lean}.
\end{small}

\end{thebibliography}
\end{document}
